\section{Limitaciones y adaptación de los LLM: datos, memoria, actualización e interpretabilidad}

Cuanto más amplias son las aplicaciones, más evidentes se vuelven las carencias del modelo. Esta sección sigue el orden de la clase y aborda cinco grupos de problemas: costo de entrenamiento y eficiencia de datos; memoria y personalización; datos sintéticos y memorization; aprendizaje continuo y model editing; self-reflection y hallucination. Todos responden a la misma pregunta: una vez congelados los parámetros, ¿cómo sigue aprendiendo, actualizándose y corrigiéndose el sistema?

\subsection{Datos, cómputo y concentración de la investigación}

Entrenar modelos grandes requiere enormes cantidades de tokens, accelerator time, energía e infraestructura de ingeniería. La clase vincula este punto con la democratización de la investigación: cuando solo unas cuantas instituciones pueden costear experimentos reproducibles, se restringen la comparación de métodos, la auditoría abierta y la participación de los estudiantes. Esta subsección revisa primero la diapositiva de costos y después pasa a BabyLM y a los modelos pequeños, para responder si «la mejora de capacidades solo puede comprarse con más cómputo».

\lecturefigure{slide-053.jpg}{Demanda de datos, cómputo, financiamiento y energía en el entrenamiento de LLM}{slide deck oficial de CS25 V4, página 53}

\begin{importantbox}{Escala aproximada del cómputo de entrenamiento}
Para el preentrenamiento de un Dense Transformer suele usarse la aproximación
\[
C\approx 6ND,
\]
donde $N$ es el número de parámetros, $D$ es el número de tokens de entrenamiento y $C$ es el orden de magnitud de las operaciones de punto flotante. La constante varía según la architecture y la implementación, pero la aproximación revela dos hechos: tanto los parámetros como los datos elevan linealmente el cómputo de entrenamiento; y tokens de mayor calidad, activación dispersa y kernels eficientes pueden cambiar el costo real.
\end{importantbox}

\begin{warningbox}{El costo no es solo una factura de GPU}
La recolección y el filtrado de datos, los experimentos fallidos, los checkpoints, la comunicación de red, el servicio de inferencia, la evaluación y la retroalimentación humana consumen recursos. Reportar solo los FLOPs del entrenamiento final subestima la cadena completa de investigación y desarrollo, y tampoco indica el costo por unidad de capacidad.
\end{warningbox}

\subsection{BabyLM: qué pregunta plantea la eficiencia de datos humana}

La diapositiva de BabyLM compara a los niños con los LLM: los niños están expuestos a mucho menos lenguaje, pero se apoyan en la percepción, la interacción, la retroalimentación social y la experiencia corporizada para formar conceptos sólidos. Esta analogía no demuestra que las máquinas deban replicar el aprendizaje infantil, pero sí expone el problema de eficiencia de datos que tiene el escalamiento basado solo en texto. Al leer las dos diapositivas en paralelo, conviene atender a la supervision density y a la interaction, y no solo comparar el número de palabras.

\lecturefigure{slide-054.jpg}{Diferencias en el entorno de datos entre el aprendizaje infantil y el entrenamiento de LLM}{slide deck oficial de CS25 V4, página 54}

\lecturefigure{slide-055.jpg}{Volumen de datos, proceso de desarrollo y experiencia multimodal en la discusión sobre BabyLM}{slide deck oficial de CS25 V4, página 55}

\begin{knowledgebox}{Lectura de la figura: los niños no son «modelos de lenguaje con pocos datos»}
Los niños cuentan con visión, acción, atención conjunta, retroalimentación emocional, intervención causal y un entorno persistente. Comparar directamente el número de palabras con los tokens de un LLM pasa por alto la supervision density y la interaction. El valor de BabyLM está en obligar a los investigadores a explorar mejores data curriculum, architecture e inductive bias.
\end{knowledgebox}

\teachervoice{La clase no presentó BabyLM como «los niños demuestran que el scaling está equivocado», sino como una pregunta de investigación: si los humanos pueden aprender con menos datos lingüísticos, ¿podrían los modelos alcanzar una mayor sample efficiency mediante datos de mayor calidad, estructuras distintas o interacción?}

\subsection{Modelos pequeños y on-device model}

Los Minified LLM y los on-device model cambian el objetivo de «la máxima capacidad general» a «costo, latencia y privacidad con una capacidad aceptable». Los modelos más pequeños pueden ganar utilidad práctica mediante destilación, quantization, pruning, datos especializados y retrieval.

\lecturefigure{slide-056.jpg}{Motivaciones y enfoques de la miniaturización y de los LLM en el dispositivo}{slide deck oficial de CS25 V4, página 56}

\begin{knowledgebox}{Glosario: cuatro enfoques de compresión de modelos}
\begin{tabularx}{\textwidth}{L{0.20\textwidth} X}
\toprule
Método & Mecanismo central \\
\midrule
Distillation & Entrena un student más pequeño con las salidas, los logits o los reasoning traces del teacher. \\
Quantization & Representa weights/activations con un bit-width menor, lo que reduce memoria y cómputo. \\
Pruning & Elimina conexiones, channels, heads o blocks de baja contribución. \\
Specialization & Usa datos de dominio, adapters, tools o retrieval para acotar el rango de capacidades que deben quedar parametrizadas. \\
\bottomrule
\end{tabularx}
\end{knowledgebox}

\subsection{Memory augmentation y personalization}

Los parámetros de un LLM congelado no se actualizan automáticamente con cada conversación. La memory augmentation (aumento de memoria) coloca parte de la información en un store externo y la recupera hacia el context cuando hace falta; la personalization (personalización) permite que el sistema aproveche de forma continua las preferencias, el historial y las restricciones del usuario. Las dos diapositivas siguientes ponen lado a lado «modificar la entrada, modificar la memoria externa y modificar los parámetros»; al leerlas, primero hay que identificar qué objeto cambia cada enfoque.

\lecturefigure{slide-057.jpg}{Contradicción entre el conocimiento congelado y la necesidad de personalización a largo plazo}{slide deck oficial de CS25 V4, página 57}

\lecturefigure{slide-058.jpg}{Enfoques de adaptación: Fine-tuning, prompt, memory bank y RAG}{slide deck oficial de CS25 V4, página 58}

\begin{importantbox}{Los cuatro mecanismos de adaptación no son intercambiables}
\begin{tabularx}{\textwidth}{L{0.19\textwidth} L{0.24\textwidth} X}
\toprule
Mecanismo & Qué cambia & Problemas adecuados \\
\midrule
Prompt/context & La entrada de esta ocasión & Instrucciones de tareas temporales y poca evidencia. \\
Retrieval/RAG & Memory externa actualizable & Hechos, documentos, historial del usuario y evidencia citable. \\
Adapter/fine-tuning & Parte o la totalidad de los parámetros & Comportamiento estable, formato, distribución de dominio y adaptación de habilidades. \\
Model editing & Parámetros locales previamente localizados & Modificación dirigida de unos pocos hechos o asociaciones; requiere verificar efectos secundarios. \\
\bottomrule
\end{tabularx}
\end{importantbox}

\begin{importantbox}{Retrieval-augmented generation en su forma mínima}
Dada una query $q$, una colección de documentos $\mathcal{D}$, un recuperador $R$ y un generador $G$:
\[
z_{1:k}=R(q,\mathcal{D}),\qquad y=G(q,z_{1:k}).
\]
$z_{1:k}$ es la evidence top-$k$. La recuperación amplía el conocimiento accesible, pero no garantiza que la evidence sea correcta, que el orden sea razonable ni que el generator realmente la utilice.
\end{importantbox}

\begin{lstlisting}[caption={Flujo mínimo de recuperación desde memoria externa y generación}]
def answer(query, memory, encoder, generator, top_k=5):
    query_vector = encoder(query)
    evidence = memory.search(query_vector, top_k=top_k)
    prompt = build_prompt(query=query, evidence=evidence)
    return generator(prompt), evidence
\end{lstlisting}

\teachervoice{El ponente distinguió expresamente entre la memory externa y la skill interna: RAG sirve para complementar hechos e historial del usuario, pero no hace que el modelo adquiera automáticamente nuevas capacidades básicas. Si el datastore es de mala calidad, la recuperación solo llevará los errores más rápido al context.}

\subsection{Synthetic data y Phi: ¿puede la calidad de los datos sustituir parte de la escala?}

Los datos sintéticos (synthetic data) se generan con un modelo potente, que produce instructions, answers o reasoning traces, con los que después se entrenan otros modelos. Reducen el costo de la recolección manual, pero también pueden copiar los errores, los sesgos de estilo y los puntos ciegos del teacher. Esta subsección toma Phi como caso representativo del enfoque centrado en la calidad de los datos: primero revisa la generación y el filtrado de datos, y después el alcance de los resultados con modelos pequeños.

\lecturefigure{slide-059.jpg}{Datos sintéticos de preentrenamiento, filtrado y distillation}{slide deck oficial de CS25 V4, página 59}

\begin{warningbox}{Tres riesgos de ciclo cerrado de los synthetic data}
\begin{enumerate}
  \item \textbf{Error inheritance}: el student aprende de nuevo los errores del teacher;
  \item \textbf{Diversity collapse}: la distribución generada es más estrecha que el mundo real y, tras entrenamientos repetidos, los patrones convergen;
  \item \textbf{Evaluation leakage}: los prompts o las respuestas coinciden con el benchmark, lo que infla las puntuaciones.
\end{enumerate}
Por ello se necesitan provenance, deduplication, quality filter y comparación con datos reales.
\end{warningbox}

\lecturefigure{slide-060.jpg}{Resultados presentados en clase de Phi-2, que entrena un modelo pequeño con datos de alta calidad y sintéticos}{slide deck oficial de CS25 V4, página 60}

\teachervoice{El ponente situó el takeaway principal de Phi en la data source y la quality, en lugar de presentar el bajo número de parámetros como algo mágico. Que un modelo con menos parámetros conserve sus capacidades depende de la cobertura de tareas, del filtrado de datos y del alcance de la evaluación.}

\begin{knowledgebox}{Lectura de la figura: Phi representa el enfoque de que «cada token valga más»}
La clase destaca los textbook-quality data y los synthetic reasoning data, que permiten a un modelo más pequeño acercarse a modelos más grandes en algunas tareas. La figura no demuestra que los modelos pequeños sean equivalentes en todas las capacidades; muestra que, además de la escala de parámetros, la data selection y el curriculum pueden cambiar los resultados de manera significativa.
\end{knowledgebox}

\subsection{¿Learning o memorizing?}

Cuando los datos de entrenamiento y los benchmarks no son transparentes, la salida del modelo puede provenir de una generalización composicional o de una memorización aproximada. La contaminación de prueba (test contamination) hace que el modelo haya visto las preguntas o contenido muy similar antes de la evaluación, de modo que la puntuación deja de representar la generalización a tareas nuevas. La diapositiva del debate que sigue debe leerse como un problema de evaluation design, no como una exigencia de elegir con prisa entre «piensa» y «solo recita».

\lecturefigure{slide-061.jpg}{Conocimiento nuevo, combinación de patrones y el debate sobre la memorization en los LLM}{slide deck oficial de CS25 V4, página 61}

\begin{importantbox}{Tres niveles de una memorization audit}
\begin{enumerate}
  \item \textbf{Exact match}: si el modelo puede reproducir fragmentos largos del texto de entrenamiento;
  \item \textbf{Near duplicate}: si el conjunto de entrenamiento contiene ligeras reformulaciones de las preguntas, las plantillas o las respuestas;
  \item \textbf{Behavioral generalization}: si el mecanismo sigue funcionando al cambiar entidades, estructura o restricciones.
\end{enumerate}
Quedarse en el primer nivel deja pasar mucha contamination, y fijarse solo en la exactitud final no permite distinguir entre memorización y generalización.
\end{importantbox}

\begin{warningbox}{No confundir la controversia de derechos de autor con la controversia cognitiva}
Si el modelo «entiende» es una cuestión cognitiva y de comportamiento; si los datos de entrenamiento cuentan con autorización y si la salida reproduce en exceso la obra original son cuestiones legales y de gobernanza. Ambas están relacionadas, pero una conclusión no puede sustituir a la otra.
\end{warningbox}

\subsection{Continual learning: no es otro nombre del reentrenamiento periódico}

El continual learning (aprendizaje continuo) exige que el modelo se actualice cuando llegan nuevas experiencias, conservando a la vez sus capacidades previas y controlando el olvido catastrófico. La clase toma como contraste el aprendizaje cotidiano de las personas y critica que la trace distillation o el fine-tuning periódico se presenten directamente como automejora permanente. Esta subsección define primero el stability--plasticity tradeoff y después explica por qué mejorar la puntuación en una sola tarea nueva no basta para demostrar aprendizaje continuo.

\lecturefigure{slide-062.jpg}{Aprendizaje continuo, adaptación en línea y olvido catastrófico}{slide deck oficial de CS25 V4, página 62}

\begin{importantbox}{Continual-learning objective}
Para una secuencia de tareas $\mathcal{T}_1,\ldots,\mathcal{T}_m$, el modelo actualizado debe reducir la pérdida en la tarea nueva y, al mismo tiempo, limitar la degradación en las tareas anteriores:
\[
\min_\theta\ \mathcal{L}_{\mathcal{T}_m}(\theta)+
\lambda\sum_{i<m}\Omega_i(\theta,\theta_{i}^{*}).
\]
$\Omega_i$ es la regularization, el replay o la parameter constraint que preserva el conocimiento previo, y $\lambda$ controla el equilibrio entre plasticity y stability. Lo verdaderamente difícil es verificar el olvido cuando no se dispone de todos los datos anteriores.
\end{importantbox}

\teachervoice{El ponente ilustró el problema con un contraste intuitivo: las personas no se sientan cada dos meses a «volver a leer todo internet» para seguir aprendiendo. Destilar nuevos traces en el modelo se parece más a un reentrenamiento y no equivale a un aprendizaje continuo en línea y a lo largo de toda la vida.}

\subsection{Interpretabilidad y model editing}

La interpretabilidad busca entender cómo las representaciones y los cálculos internos producen la salida; la mechanistic interpretability se centra más en components, circuits y causal intervention concretos. El model editing, por su parte, pretende cambiar de forma dirigida una asociación factual concreta sin reentrenar todo el modelo.

\lecturefigure{slide-063.jpg}{La caja negra de los modelos grandes, el control y la mechanistic interpretability}{slide deck oficial de CS25 V4, página 63}

\begin{knowledgebox}{No mezclar tres tipos de «explicación»}
\begin{tabularx}{\textwidth}{L{0.24\textwidth} X}
\toprule
Nivel & Pregunta que responde \\
\midrule
Behavioral explanation & ¿Cómo influyen los cambios en la entrada sobre la salida y en qué distribuciones falla el modelo? \\
Attribution/probing & ¿Qué tokens, activations o features se relacionan con la salida? \\
Mechanistic/causal analysis & ¿Intervenir un component cambia de forma estable el comportamiento objetivo? \\
\bottomrule
\end{tabularx}
La correlación ayuda a localizar, pero no constituye por sí misma un mecanismo causal.
\end{knowledgebox}

\lecturefigure{slide-064.jpg}{Esquema de clase de ROME y del model editing factual}{slide deck oficial de CS25 V4, página 64}

\begin{importantbox}{Forma abstracta de un rank-one edit}
Si se quiere actualizar de forma dirigida la matriz de pesos $W$ de cierta capa, puede escribirse
\[
W' = W + uv^\top,
\]
donde $u,v$ son vectores obtenidos a partir de la key/value relation objetivo. Un rank-one update tiene bajo costo, pero su validación debe revisar a la vez edit success, paraphrase generalization, locality y unrelated knowledge damage.
\end{importantbox}

\begin{warningbox}{Corregir un hecho no significa que la «base de conocimiento» del modelo quede reparada}
Un hecho puede estar distribuido en varias capas y contextos; una edición local también puede alterar entidades cercanas o patrones lingüísticos. El model editing requiere una regression suite, no solo mostrar que el prompt objetivo ya produce la nueva respuesta.
\end{warningbox}

\subsection{MoE, modularity y autorreflexión}

El curso vuelve a MoE y plantea una pregunta más audaz: ¿podría un solo modelo formar módulos reutilizables en su interior, como las áreas funcionales del cerebro humano? Esta analogía sirve para estimular el diseño de architecture, pero no puede ocultar las dificultades reales del router, la expert specialization y la global coordination.

\lecturefigure{slide-065.jpg}{Estructura de routed experts de MoE y activación dispersa}{slide deck oficial de CS25 V4, página 65}

\lecturefigure{slide-066.jpg}{Propuesta abierta: de múltiples experts a la modularización dentro de un solo modelo}{slide deck oficial de CS25 V4, página 66}

\begin{warningbox}{La modularización necesita interfaces e indicadores de validación}
Si entre los módulos no hay entradas y salidas definidas, routing criterion, capacity y failure isolation, la «modularización» es solo un diagrama conceptual. Entre los indicadores medibles están la expert load, el token drop, la routing entropy, la specialization, el costo de comunicación y la propagación de fallas.
\end{warningbox}

La self-reflection (autorreflexión) hace que el modelo lea su propia respuesta, genere una critique y la modifique de forma iterativa. Puede aumentar el test-time compute y también sacar a la luz errores; pero si el critique model comparte puntos ciegos con el modelo original, el ciclo puede producir solo errores más largos y más seguros de sí mismos.

\lecturefigure{slide-067.jpg}{La self-reflection itera la salida mediante critique y revision}{slide deck oficial de CS25 V4, página 67}

\lecturefigure{slide-068.jpg}{Cambios de desempeño y límites de la self-reflection de varios niveles}{slide deck oficial de CS25 V4, página 68}

\begin{importantbox}{Reflection loop}
Sea la respuesta inicial $y_0=G(x)$; en la ronda $t$ se genera la critique $c_t=C(x,y_t)$ y después se revisa
\[
y_{t+1}=R(x,y_t,c_t).
\]
$G$ es el generador, $C$ es el proceso de critique y $R$ es el proceso de revision. El ciclo necesita una stopping rule externa y un verifier; de lo contrario, «seguir reflexionando» puede aumentar el costo sin aumentar la exactitud.
\end{importantbox}

\subsection{Hallucination, calibration y evidencia externa}

La diapositiva de hallucination describe el problema como «el modelo no sabe que no sabe». Dicho con más precisión, el modelo puede seguir produciendo lenguaje de alta confianza aunque carezca de evidencia. La calibration (calibración) se ocupa de si la confianza coincide con la tasa real de aciertos; retrieval y verifier intentan, por su parte, incorporar evidencia externa.

\lecturefigure{slide-069.jpg}{Hallucination, verificación de hechos, calibration y el enfoque de retrieval}{slide deck oficial de CS25 V4, página 69}

\begin{importantbox}{Definición de calibration}
Si el modelo asigna una confianza $q$ a un conjunto de respuestas, la calibración ideal exige
\[
P(\text{correct}\mid \text{confidence}=q)\approx q.
\]
Estar bien calibrado no equivale a tener alta exactitud; un modelo que responde «no estoy seguro» con frecuencia puede estar muy bien calibrado. A la inversa, un modelo de alta exactitud puede mostrar exceso de confianza en sus pocos errores.
\end{importantbox}

\begin{warningbox}{La temperature no es una perilla de factualidad}
Reducir la sampling temperature hace que la distribución sea más aguda y más cercana a la greedy decoding, pero no puede convertir un conocimiento erróneo en uno correcto. La factualidad requiere el apoyo conjunto de evidence, retrieval, verification, datos de entrenamiento y restricciones de la tarea.
\end{warningbox}

\subsection{Resumen de la sección}

La principal carencia de los LLM no es un bug aislado, sino un conjunto de problemas acoplados entre sí: los datos y el cómputo determinan el alcance de lo que puede aprenderse, la memory externa ofrece un canal de actualización, el aprendizaje continuo y el model editing modifican los parámetros, la interpretabilidad ayuda a localizar mecanismos, y la reflection y la calibration cambian el comportamiento en el momento de la inferencia. Ningún método por sí solo resuelve automáticamente a la vez el conocimiento, las habilidades, la factualidad y la adaptación a largo plazo.
